% Global context: archival paper for the measured MiniBug-RL experiment.
% Sources are cited inline and listed in the bibliography at the end of this file.
\documentclass[11pt]{article}

\usepackage[T1]{fontenc}
\usepackage[utf8]{inputenc}
\usepackage{amsmath}
\usepackage{booktabs}
\usepackage{geometry}
\usepackage{hyperref}
\usepackage{xcolor}

\geometry{margin=1in}
\hypersetup{
  colorlinks=true,
  linkcolor=blue!55!black,
  citecolor=blue!55!black,
  urlcolor=blue!55!black,
  pdftitle={MiniBug-RL: Tiny Reproducible Reinforcement Learning for Python Repair},
  pdfauthor={BurnyCoder}
}
\setlength{\parindent}{0pt}
\setlength{\parskip}{0.55em}
% Global context: checked numerical macros keep the paper synchronized with public evidence.
% Source: reports/evidence/run-summary.json and run 20260904T021845Z-minibug-rl-main.
\newcommand{\BaseRevision}{ea3f2471cf1b1f0db85067f1ef93848e38e88c25}
\newcommand{\SourceCommit}{7027bc55baecc00fad51cbe2b8f030dca2c91c1e}
\newcommand{\HubRevision}{5b6e22a4c6c01bec95d10e93a0fc78666eb9c543}
\newcommand{\DockerImageDigest}{a869cd1dffb8c87afad1bb1302106cb9f5cb580641c7391bb73f4ab077f140d9}
\newcommand{\TrainSteps}{100}
\newcommand{\TrainLoss}{0.017962}
\newcommand{\TrainSeconds}{1,174.08}
\newcommand{\TrainPeakBytes}{1,849,278,976}
\newcommand{\TrainPeakGiB}{1.72}
\newcommand{\TrainableParameters}{8,798,208}
\newcommand{\TrainablePercent}{1.7497\%}
\newcommand{\ChangedTensors}{336}
\newcommand{\AdapterDelta}{0.281200}
\newcommand{\BaseValSolved}{7/12}
\newcommand{\ModelValSolved}{8/12}
\newcommand{\BaseValHidden}{0.7708}
\newcommand{\ModelValHidden}{0.8333}
\newcommand{\ValDifference}{+0.0625}
\newcommand{\ValInterval}{[-0.1458,\ 0.2917]}
\newcommand{\BaseTestSolved}{5/12}
\newcommand{\ModelTestSolved}{6/12}
\newcommand{\BaseTestHidden}{0.6083}
\newcommand{\ModelTestHidden}{0.6042}
\newcommand{\TestDifference}{-0.0042}
\newcommand{\TestInterval}{[-0.2917,\ 0.2625]}
\newcommand{\BaseSampledSolve}{10/12}
\newcommand{\ModelSampledSolve}{10/12}
\newcommand{\BaseSampledHidden}{0.4635}
\newcommand{\ModelSampledHidden}{0.6281}
\newcommand{\BaseInvalid}{19/60}
\newcommand{\ModelInvalid}{1/60}
\newcommand{\BaseExternalPassed}{37/164}
\newcommand{\ModelExternalPassed}{38/164}
\newcommand{\BaseExternalRate}{0.2256}
\newcommand{\ModelExternalRate}{0.2317}
\newcommand{\ExternalDifference}{+0.0061}
\newcommand{\ExternalInterval}{[-0.0183,\ 0.0366]}
\newcommand{\BaseExternalTimeouts}{3}
\newcommand{\ModelExternalTimeouts}{4}


\title{MiniBug-RL: Tiny, Reproducible Reinforcement Learning\\for Verifiable Python Repair}
\author{BurnyCoder\\\small\url{https://github.com/BurnyCoder/llm-rl-software-engineering}}
\date{September 4, 2026}

\begin{document}
\maketitle

\begin{abstract}
We ask whether a real reinforcement-learning-from-verifiable-rewards loop for software repair can be made reproducible on one 8~GB consumer GPU. We train a rank-16 LoRA adapter on the 0.49B-parameter Qwen2.5-Coder model for \TrainSteps{} Group Relative Policy Optimization steps. Four candidate repairs per prompt receive rewards from hidden unit tests executed in fresh resource-limited Docker containers. Training completed in \TrainSeconds{} seconds with \TrainPeakGiB{}~GiB peak allocated GPU memory and changed all \ChangedTensors{} trainable adapter tensors. Validation improved from \BaseValSolved{} to \ModelValSolved{} greedily solved tasks and selected the 100-step checkpoint. On the untouched 12-task internal test, full greedy solves rose from \BaseTestSolved{} to \ModelTestSolved{}, but mean greedy hidden-test correctness changed from \BaseTestHidden{} to \ModelTestHidden{} (paired difference \TestDifference{}, 95\% bootstrap interval \TestInterval{}). On frozen HumanEvalFixDocs-Python, greedy pass@1 changed from \BaseExternalPassed{} to \ModelExternalPassed{} (difference \ExternalDifference{}, interval \ExternalInterval{}). Invalid internal candidates fell from \BaseInvalid{} to \ModelInvalid{}. The result establishes an end-to-end, publicly reloadable RL pipeline and a clear output-validity improvement; it does not establish a statistically conclusive general repair-capability gain.
\end{abstract}

\section{Research question and contribution}

Software is unusually suitable for verifiable reinforcement learning because candidate behavior can be checked by executable tests rather than a learned preference model. The practical obstacle is that serious repository-scale agent benchmarks often combine large policies, long contexts, tool orchestration, and substantial compute. This experiment narrows the domain to deterministic single-function Python repair and asks:

\begin{quote}
Can a sub-billion-parameter code model receive genuine group-relative policy updates from hidden tests, on one 8~GB GPU, while retaining defensible split discipline, execution isolation, complete evidence, and a resulting model that another user can load?
\end{quote}

The hypothesis was that 100 low-rank GRPO steps would improve validation hidden-test correctness by at least 0.05 or solve one additional validation task, without worsening the aggregate format/runtime/policy status rate. The experiment contributes (i) a fully pinned implementation; (ii) an original 36/12/12 train/validation/test curriculum with structural leakage checks; (iii) a Docker-mediated reward path; (iv) locked internal and external paired evaluation; and (v) public merged and adapter artifacts.

The scope is intentionally modest. A function-repair policy is not a repository agent, and passing small unit-test suites is not evidence of broad software-engineering autonomy.

\section{Foundations and implementation choice}

GRPO was introduced as a policy-optimization variant that forms relative advantages from groups of sampled completions while avoiding a learned value model \cite{deepseekmath}. TRL exposes the algorithm, custom reward functions, PEFT integration, and DR-GRPO loss in a maintained trainer interface \cite{trl}. LoRA represents the update with small trainable low-rank matrices while leaving base weights frozen \cite{lora,peft}. These properties make the combination suitable for limited memory.

Qwen2.5-Coder provides instruction-tuned code models down to 0.5B parameters \cite{qwen}. The pinned base is
\begin{center}
\small\texttt{Qwen/Qwen2.5-Coder-0.5B-Instruct}\\
\texttt{commit \BaseRevision}.
\end{center}
The public DebugArena model card motivated the particular 0.5B model, four-completion group, and $10^{-5}$ learning rate \cite{debugarena}. Its illustrative snippet does not define its reward function, dataset construction, adapter attachment, isolation, held-out protocol, or released trained checkpoint. MiniBug-RL therefore uses the pinned official TRL, PEFT, Transformers, and Hub APIs rather than treating that snippet as a complete recipe.

\section{Method}

\subsection{Curriculum and prompt contract}

The original curriculum contains 60 short, import-free functions: 36 training, 12 validation, and 12 final-test tasks. Families include incorrect comparisons, boundary and off-by-one errors, bad state updates, accumulation errors, indexing mistakes, and Boolean logic. Each task has a natural-language specification, buggy source, two public JSON examples, and four to eight hidden JSON cases. It contains no gold repair.

The model receives only the specification, buggy function, and public examples. It is instructed to return exactly one complete replacement function. A versioned manifest hashes canonical task records and normalized abstract syntax trees. Bound names are alpha-normalized and literal values are masked by type before structural hashing; no normalized template crosses a split. Every original bug was run in the real container and shown to fail at least one hidden case before training.

\subsection{Reward and isolation}

For a valid candidate with $n$ hidden tests, let $p$ be the fraction passed and $I$ indicate a full solve. The ordinary successful-execution reward is
\begin{equation}
  r = p + I[p=1] + 0.1,
\end{equation}
so a syntactically valid but wholly wrong function earns 0.1 and a complete repair earns 2.1. Invalid structure returns $-0.3$, a prohibited construct returns $-1.0$, and a valid function that crashes or times out receives $0.1-0.25=-0.15$. Expected values remain in the trusted host process; only candidate source and JSON call inputs cross the container boundary. Actual results return to the host for comparison.

Every rollout runs in a newly named container from this immutable image ID:
\begin{center}
\small\texttt{sha256:\DockerImageDigest}.
\end{center}
Networking is disabled (\texttt{--network none}) and the root filesystem is read only. The runtime drops all capabilities, sets \texttt{no-new-privileges}, uses non-root UID/GID 65532, and allows no mounts. It also enforces a 32-PID limit, 128~MiB memory with no swap, 0.5 CPU, a 16~MiB no-execute temporary filesystem, and a host deadline. Docker documents these controls and the limitations of container isolation \cite{dockerrun,dockernone}. AST filtering and reduced builtins are defense in depth, not a security boundary.

\subsection{Model update}

The implementation uses BF16 Transformers/PEFT rather than 4-bit loading. LoRA rank is 16, alpha is 32, dropout is zero, bias is disabled, and \texttt{all-linear} targets linear and Conv1D layers except the output layer. This produces \TrainableParameters{} trainable parameters (\TrainablePercent{} of 502,830,976 total).

The main GRPO configuration is shown in Table~\ref{tab:training}. The effective optimizer batch is $1\times4=4$, divisible by the four generations. Reward scaling is disabled, $\beta=0$ omits a reference-model KL term, the DR-GRPO loss uses the fixed completion-length divisor, and fully truncated completions are masked. Gradient checkpointing and SDPA reduce memory. A callback would stop if more than 80\% of groups had zero reward variance for 20 consecutive logs.

\begin{table}[ht]
\centering
\caption{Checked main training configuration.}
\label{tab:training}
\begin{tabular}{@{}ll@{}}
\toprule
Setting & Value \\
\midrule
Steps / seed & 100 / 42 \\
Generations per prompt & 4 \\
Per-device batch / accumulation & 1 / 4 \\
Prompt / completion limit & 512 / 256 tokens \\
Learning rate / warmup & $10^{-5}$ / 0.05 \\
Sampling & temperature 0.9, top-$p$ 0.95 \\
Gradient norm / precision & 0.1 / BF16 \\
Loss / KL coefficient & DR-GRPO / $\beta=0$ \\
LoRA rank / alpha / dropout & 16 / 32 / 0 \\
Checkpoint interval & 25 steps; retain latest two \\
\bottomrule
\end{tabular}
\end{table}

\subsection{Pre-registered selection and evaluation}

The untouched base, two-step smoke adapter, and 100-step adapter use identical validation prompts, task-level seeds, decoding, and sandbox. Candidate adapters are ranked lexicographically by mean greedy hidden-test fraction and then greedy full-solve rate. The learning gate requires a hidden-fraction point gain of at least 0.05 or one additional solved task, and an aggregate rate of invalid, timeout, runtime, and policy statuses no worse than the base.

The optional 200-step continuation decision was made before opening final data. The 100-step adapter passed the gate, while mean training rewards in steps 51--75 and 76--100 were 1.5230 and 1.5355. The checkpoint was locked rather than extending the budget opportunistically.

Internal evaluation uses one greedy and four sampled candidates per task. ``Sampled success@4'' is the observed fraction with at least one full solve among exactly four draws, not the unbiased pass@$k$ estimator discussed by Chen et al. \cite{codexeval}. Confidence intervals use 10,000 paired task-level bootstrap resamples with seed 42.

Only after selection, the same checkpoint is evaluated on HumanEvalFixDocs-Python from HumanEvalPack \cite{octopack}, pinned to dataset revision
\begin{center}
\small\texttt{9a41762f73a8cb23bb5811b73d5aab164efcf378}.
\end{center}
There is one greedy candidate per each of 164 tasks. The prompt conversion and first-block stopping follow the pinned BigCode harness. MiniBug-RL applies a 3-second hardened-container deadline, stricter than the harness's 10-second Python limit; the score is therefore not leaderboard-identical.

\section{Results}

\subsection{Training viability}

The two-step smoke changed 168 adapter tensors and survived save/reload. Main training reached exactly \TrainSteps{} steps with finite loss \TrainLoss{}, changed all \ChangedTensors{} trainable tensors, and had adapter $\ell_2$ delta \AdapterDelta{}. Peak allocated GPU memory was \TrainPeakBytes{} bytes (\TrainPeakGiB{}~GiB), well below the recorded 8,151~MiB device capacity. The longest zero-reward-variance streak was eight, below the 20-log stop threshold. Four of 400 rollouts were clipped at the completion limit, all by step 31. No training, Trackio, CUDA, or reward-collapse failure occurred.

\subsection{Validation and locked internal test}

\begin{table}[ht]
\centering
\caption{Internal results. Hidden fraction is a macro-average over tasks.}
\label{tab:internal}
\begin{tabular}{@{}lrr@{}}
\toprule
Metric & Base & Selected adapter \\
\midrule
Validation greedy solves & \BaseValSolved{} & \ModelValSolved{} \\
Validation greedy hidden fraction & \BaseValHidden{} & \ModelValHidden{} \\
Final greedy solves & \BaseTestSolved{} & \ModelTestSolved{} \\
Final greedy hidden fraction & \BaseTestHidden{} & \ModelTestHidden{} \\
Final observed sampled success@4 & \BaseSampledSolve{} & \ModelSampledSolve{} \\
Final sampled hidden fraction & \BaseSampledHidden{} & \ModelSampledHidden{} \\
Final invalid candidates & \BaseInvalid{} & \ModelInvalid{} \\
\bottomrule
\end{tabular}
\end{table}

Validation selected the main adapter with a greedy hidden-fraction difference of \ValDifference{} (95\% interval \ValInterval{}) and one additional solved task. The interval crosses zero; the deterministic pre-registered gate passed, but this is not statistical significance.

On final tasks, greedy full solves rose by one, while mean greedy hidden fraction was essentially unchanged: difference \TestDifference{}, interval \TestInterval{}. Both policies passed 30 of 49 raw greedy hidden cases; the small difference in macro means arises because tasks are weighted equally and one task has five cases. Sampled hidden correctness increased, and invalid formatting fell markedly, although runtime errors rose from four to seven and the selected set added one timeout and one policy violation. Thus output-structure reliability improved overall but not every reliability category did.

\subsection{Frozen external evaluation}

\begin{table}[ht]
\centering
\caption{HumanEvalFixDocs-Python, greedy $n=1$, 164 tasks, 3-second MiniBug-RL sandbox.}
\label{tab:external}
\begin{tabular}{@{}lrr@{}}
\toprule
Metric & Base & Selected adapter \\
\midrule
Passed & \BaseExternalPassed{} & \ModelExternalPassed{} \\
Pass@1 & \BaseExternalRate{} & \ModelExternalRate{} \\
Functional failures & 124 & 122 \\
Timeouts & \BaseExternalTimeouts{} & \ModelExternalTimeouts{} \\
\bottomrule
\end{tabular}
\end{table}

The paired external point difference is \ExternalDifference{} with interval \ExternalInterval{}. The adapter gained Python/69, Python/78, and Python/98, and lost Python/62 and Python/86; 159 task outcomes tied. External timeouts increased from \BaseExternalTimeouts{} to \ModelExternalTimeouts{}. This one-task net gain is compatible with no underlying effect.

\section{Error analysis and interpretation}

The clearest learned behavioral change is format adherence on the internal final protocol: invalid candidates dropped from 19 to one among 60 greedy-plus-sampled candidates. This matters operationally because a function-repair policy cannot earn execution feedback from prose, multiple functions, or malformed source. It may also explain the increase in sampled partial correctness.

Correctness changes were heterogeneous. Greedy internal gains occurred on Boolean parsing, nearest multiples, wrapped indexing, and pair sums. Regressions occurred on Hamming distance, alternating strings, and trace computation. The adapter sometimes transformed an invalid output into executable but incorrect source, moving errors from parsing into runtime behavior. Longer external completions also made the adapter evaluation nearly twice as slow and produced one additional timeout.

Accordingly, the evidence supports the narrow claim that verifiable GRPO was implemented successfully and altered the policy in a useful structural direction. It does not support a broad claim of improved software engineering. All three paired intervals---validation, internal final, and external---include zero.

\section{Experiment history and corrective loop}

An initial diagnostic run completed preflight, preparation, baseline, and a two-step smoke check, but exposed an observability defect: TRL emitted undefined clip extrema for an entirely clipped batch, and Trackio 0.37 dropped that complete metric record. The run was stopped manually before main training. The fix subclasses the official callback, removes only \texttt{None} fields, and retains zero and negative metrics. Curriculum review at the same boundary found two hidden suites that did not catch their authored bugs and one cross-split structural collision; all were corrected and re-audited before the authoritative run.

This diagnostic evidence is not pooled with the final experiment because its source, data, and manifest hashes differ. The authoritative loop was: define a small verifiable domain; state the gate; audit data and isolation; measure the base; prove gradient/save/reload viability; train 100 steps; decide the budget from training plus validation; lock; evaluate once; export; publish; re-download; and report both positive and neutral outcomes.

\section{Threats to validity}

\textbf{Statistical power.} The internal sets contain only 12 tasks each and the experiment has one training seed. One solved task moves a pass rate by 8.33 percentage points. Bootstrap intervals quantify paired task uncertainty but not variation across training seeds or curriculum construction.

\textbf{Selection.} Validation is intentionally optimistic because it selects the checkpoint and budget. Only final-test and external results estimate post-selection behavior.

\textbf{Benchmark exposure.} HumanEvalPack was absent from local training, reward design, and selection, but upstream base-model pretraining exposure cannot be ruled out. Both policies share that risk.

\textbf{Protocol comparability.} The external prompt uses the raw HumanEvalFixDocs instruct form rather than Qwen's chat template, and the 3-second deadline differs from the BigCode harness. The measured numbers must not be compared directly with stock leaderboard entries.

\textbf{Security.} Containers share the host kernel and are not perfect hostile-code virtual machines. The setup uses no mounts or Docker socket, but higher-assurance deployments should add a disposable VM, rootless Docker, user namespaces, seccomp/AppArmor, or a microVM boundary.

\textbf{Scope.} Synthetic pure-function repair does not test repository navigation, dependency management, multi-file patches, tests written by the model, or long-horizon agent behavior.

\section{Reproducibility and artifacts}

The producing source commit is \texttt{\SourceCommit}. Public configuration, curriculum, and manifest SHA-256 hashes are retained with the experiment report. The run used Python 3.12, PyTorch 2.14.0+cu130, Transformers 5.16.1, TRL 1.12.0, PEFT 0.20.0, and Docker 29.4.1 on an NVIDIA GeForce RTX 5070 Laptop GPU. Every prompt/output pair, reward component, per-step metric, environment field, checkpoint, selection decision, evaluation candidate, and phase state was logged locally.

The public Hugging Face repository is
\begin{center}
\url{https://huggingface.co/BurnyCoder/qwen2.5-coder-0.5b-swe-rl}
\end{center}
at immutable revision \texttt{\HubRevision}. Its root contains the merged model, tokenizer, card, configuration, and measured results; \texttt{adapter/} contains the separate LoRA. Publication was accepted only after re-downloading that exact commit, byte-comparing both weight files, loading the root model, and producing a non-empty greedy repair.

\section{Conclusion}

MiniBug-RL answers the engineering question affirmatively: a complete software-repair RL loop, including verifiable rewards, low-rank updates, split discipline, sandboxed execution, paired evaluation, model export, and public reload verification, fits comfortably on an 8~GB GPU. The scientific result is more restrained. The selected adapter solved one additional task on each held-out benchmark and dramatically reduced malformed internal outputs, but average greedy correctness did not reliably improve and every paired interval includes zero. The appropriate next experiment is a larger structurally diverse training curriculum and multiple seeds, selected without reopening the current final sets.

\appendix
\section{Exact reproduction commands}

\begin{verbatim}
git clone https://github.com/BurnyCoder/llm-rl-software-engineering.git
cd llm-rl-software-engineering
git checkout 7027bc55baecc00fad51cbe2b8f030dca2c91c1e
uv sync --locked
cp .env.example .env
chmod 600 .env
RUN_ID="$(date -u +%Y%m%dT%H%M%SZ)-minibug-rl-main"
uv run swe-rl all --config configs/local-8gb.toml \
  --run-id "$RUN_ID" --skip-publish
\end{verbatim}

Publication is a separate, authorized action after setting \texttt{HF\_TOKEN} in the ignored mode-600 \texttt{.env}:

\begin{verbatim}
uv run swe-rl publish --config configs/local-8gb.toml \
  --run-id "$RUN_ID"
\end{verbatim}

\section{Evidence map}

Machine-readable checked summaries live in \texttt{reports/evidence/}. Run-specific hypotheses, failures, fixes, and decisions live in \texttt{reports/experiments/}. The operational procedure, stable methodology, threat model, and provenance interpretation live in separate files under \texttt{docs/}. Raw logs and model weights remain outside Git and are referenced by run ID and immutable public commits.

\begin{thebibliography}{10}

\bibitem{deepseekmath}
Z. Shao et al.
\newblock DeepSeekMath: Pushing the Limits of Mathematical Reasoning in Open Language Models.
\newblock \emph{arXiv:2402.03300}, 2024.
\newblock \url{https://arxiv.org/abs/2402.03300}

\bibitem{trl}
Hugging Face.
\newblock TRL 1.12.0 GRPO Trainer documentation.
\newblock \url{https://huggingface.co/docs/trl/v1.12.0/en/grpo_trainer}

\bibitem{lora}
E. Hu et al.
\newblock LoRA: Low-Rank Adaptation of Large Language Models.
\newblock \emph{International Conference on Learning Representations}, 2022.
\newblock \url{https://arxiv.org/abs/2106.09685}

\bibitem{peft}
Hugging Face.
\newblock PEFT LoRA API documentation.
\newblock \url{https://huggingface.co/docs/peft/main/package_reference/lora}

\bibitem{qwen}
B. Hui et al.
\newblock Qwen2.5-Coder Technical Report.
\newblock \emph{arXiv:2409.12186}, 2024.
\newblock \url{https://arxiv.org/abs/2409.12186}

\bibitem{debugarena}
B. V. Tadepalli.
\newblock DebugArena: Teaching LLMs to Fix Bugs Through Trial and Error.
\newblock Hugging Face model card, 2026.
\newblock \url{https://huggingface.co/BharathVikas/debugarena}

\bibitem{dockerrun}
Docker, Inc.
\newblock \texttt{docker container run} reference.
\newblock \url{https://docs.docker.com/reference/cli/docker/container/run/}

\bibitem{dockernone}
Docker, Inc.
\newblock None network driver.
\newblock \url{https://docs.docker.com/engine/network/drivers/none/}

\bibitem{codexeval}
M. Chen et al.
\newblock Evaluating Large Language Models Trained on Code.
\newblock \emph{arXiv:2107.03374}, 2021.
\newblock \url{https://arxiv.org/abs/2107.03374}

\bibitem{octopack}
N. Muennighoff et al.
\newblock OctoPack: Instruction Tuning Code Large Language Models.
\newblock \emph{arXiv:2308.07124}, 2023.
\newblock \url{https://arxiv.org/abs/2308.07124}

\end{thebibliography}

\end{document}
